\begingroup
\raggedright
\subsection{Hashing Toolkit: Choosing the Right Hash}
C++17 (GCC/Clang). These implementations supplement the existing
\texttt{HashedString} below. All new headers use namespace
\texttt{hashing\_ref}; keep their dependency headers in the same folder.
The shared \texttt{HashCore.hpp} appears first. Ranges are zero-based
$[l,r)$, including empty ranges. Actual string sizes and segment-tree
index arithmetic must fit \texttt{int} and available memory.
Represented lengths and multiset counts must fit \texttt{uint64\_t}.
The checked length operations throw \texttt{overflow\_error} if not.

\begin{center}
\begin{tabular}{p{0.37\columnwidth}p{0.52\columnwidth}}
\textbf{Need} & \textbf{Use} \\
\hline
Order-sensitive string equality & \texttt{Hash}, \texttt{PrefixHash} \\
Concatenated / repeated strings & \texttt{operator*}, \texttt{repeated} \\
Same elements with multiplicities & \texttt{MultisetHash}, \texttt{MultisetPrefix} \\
A dynamic set of distinct keys & \texttt{SetHash} with membership tracking \\
A board/state changed in few slots & \texttt{ZobristContext}, \texttt{ZobristState} \\
Character replacements and palindromes & \texttt{DynamicStringHash} \\
Rotation-invariant candidate filter & \texttt{RotationTraceHash}; verify matching candidates \\
\end{tabular}
\end{center}

\textbf{Shared parameters.} String hashes use the prime
$M=2^{61}-1$ and two separately sampled bases, fixed for the entire
process. Multiplication uses unsigned 128-bit intermediates.
A byte is encoded as its unsigned value plus one, in $[1,256]$;
zero bytes are supported. Each equality key includes length as well
as both residues. Hashes from different processes, seeds or contexts
are not interchangeable. Do not edit stored hash fields directly.
The randomization is for contest fingerprints, not security.
\Needspace{6\baselineskip}
\lstinputlisting[style=C++,showstringspaces=false]{../strings/HashCore.hpp}

\subsubsection{Collision Guide: What the Probability Refers To}
Unequal hashes prove inequality for the represented comparison;
equal hashes generally mean ``candidate match''. Ordinary string,
collection, and Zobrist equality methods are Monte Carlo. Verify
against the original data if an exact answer is required.
The rotation class provides such a verification method explicitly.

\textbf{Ideal uniform-hash model.} Suppose distinct objects receive
independent uniform values among $R$ possibilities. A fixed unequal
pair collides with probability $1/R$. For $Q$ specified comparisons,
the union bound gives at most $Q/R$. Among all pairs of $K$ objects,
it gives at most $\binom K2/R$, capped at one.
The exact birthday probability in this ideal model is
\[
 1-\prod_{i=0}^{K-1}\left(1-\frac{i}{R}\right)
 \approx 1-\exp\left(-\frac{K(K-1)}{2R}\right).
\]
Thus comparing a million hashes against one fixed object is a
very different risk from deduplicating a million hashes among
themselves. The number of potentially bad pairs matters.

\textbf{Polynomial hashes are not independent random labels.}
For two fixed, different strings of the same length $L$, subtracting
their polynomial hashes gives a nonzero polynomial of degree at most
$L-1$ over the prime field. A nonzero degree-$d$ polynomial has at most
$d$ roots. Here each base is uniform over $R=M-258$ choices, so a
single-base collision has probability at most
$\min(1,(L-1)/R)$. With two independently sampled bases, the fixed-pair
bound is the square of that quantity. Different lengths are
separated by the stored length. These are worst-case bounds, not
claims that every collision probability equals $1/M$ or $1/M^2$.
For many pairs, use a union bound over a fixed set of possible
comparisons; independence between collision events is unnecessary.
Input objects must be independent of the sampled parameters for
this argument. The degree bound can become weak for huge virtual
strings, even when their short descriptions fit in memory.

\textbf{Practical rules.} Use the same sampled bases for all objects
that will be compared. Keep both residues; compressing the pair to
one machine word loses information. A fixed polynomial base can be
targeted by specially constructed inputs. Plain overflow arithmetic
modulo $2^{64}$ is not a prime field, so the polynomial root argument
does not apply there. The unsigned wrapping in the multiset and
Zobrist implementations below has a different purpose and analysis.
The two seeded mixers in collection hashing are a practical
fingerprint scheme; they are not a proof of independent random
labels for every possible key. Matrix traces from one family also
must not be counted as independent hashes.

\subsection{Composable Double Hash}
For a string $s$, each component represents
$H(s)=\sum_{i=0}^{|s|-1}(\operatorname{byte}(s_i)+1)B^i\pmod M$.
A \texttt{Hash} stores two values of $H$, their $B^{|s|}$ factors,
the inverse factors, and the length.
The empty string has hash zero, powers one and length zero.

\textbf{Concatenation is a monoid operation.} If $a$ is followed by
$c$, then $H(ac)=H(a)+B^{|a|}H(c)$. Powers and inverse powers multiply.
This operation is associative, has the empty string as its identity,
and preserves order: generally $a*c\ne c*a$.
For example, concatenate the hashes of \texttt{"ab"} and
\texttt{"cd"} to get the hash of \texttt{"abcd"}, without its text.
Each concatenation costs $O(1)$.

\textbf{Removing a known prefix.} If $a$ is a prefix of $c$, then
$H(c\setminus a)=(H(c)-H(a))B^{-|a|}$.
The inverse factor is already stored, so \texttt{c/a} costs $O(1)$.
The operator assumes a real prefix; it does not check membership.
Prefix tables guarantee this precondition when extracting substrings.
\texttt{PrefixHash(s)} builds in $O(n)$ time and space;
\texttt{get(l,r)} gives a normalized substring hash in $O(1)$.

\textbf{Repetition.} Binary exponentiation works on any associative
operation with an identity. \texttt{repeated(h,k)} concatenates $k$
copies in $O(\log(k+1))$ time and $O(1)$ extra space; $k=0$ gives
empty. It checks the represented length first and avoids an unused
final squaring that could overflow even when the answer fits.
A numerical ordering of hash residues is not lexicographic string
ordering; use the comparison helper in the next subsection.
\Needspace{6\baselineskip}
\lstinputlisting[style=C++,showstringspaces=false]{../strings/ComposableHash.hpp}

\subsection{Hashing Recipes: LCP, Comparison and Periods}
\textbf{LCP.} Substring equality allows a binary search on the common
prefix length. \texttt{lcp(a,i,b,j,limit)} returns at most
\texttt{limit} matched bytes in $O(\log(n+1))$ hash queries.
\texttt{compare(a,l,r,b,i,j)} then compares the first differing byte;
if one substring is exhausted, the shorter one comes first.
It returns $-1$, zero or $1$, with the same Monte Carlo caveat as
substring equality. Byte comparisons use unsigned character values.

\textbf{Period versus a whole number of repeated blocks.}
For a nonempty substring $s[l..r)$ and $1\le k\le r-l$, period $k$
means $s[l..r-k)=s[l+k..r)$. This needs one hash comparison.
To be a power of its first $k$ bytes, also require $(r-l)\bmod k=0$.
For instance, \texttt{"abcab"} has period three but is not a whole
number of copies of \texttt{"abc"}. \texttt{has\_period} and
\texttt{is\_block\_power} make this distinction in $O(1)$.
They return false for invalid period sizes and for an empty range.

\textbf{Huge repetitions.} The first $L$ bytes of the infinite
repetition of nonempty $A$ consist of $\lfloor L/|A|\rfloor$ full
copies followed by its prefix of length $L\bmod |A|$.
\texttt{periodic\_prefix} hashes this in $O(\log(L+1))$.
\texttt{compare\_repeated(a,p,b,q)} binary-searches the LCP of
$A^p$ and $B^q$ using these hashes, in $O(\log^2(L+2))$ after prefix
preprocessing, where $L$ is the shorter represented length.
It never constructs the repeated strings. It remains probabilistic,
and very large represented lengths weaken the polynomial degree bound.

\textbf{Exact alternative for this particular comparison.}
The Fine--Wilf theorem implies that two infinite periodic strings
with periods $u=|A|$ and $v=|B|$ are identical if they agree through
$u+v-\gcd(u,v)$ positions. Therefore
\texttt{compare\_repeated\_exact} scans at most this many bytes,
or stops when one finite repetition ends. If there is no mismatch,
the finite lengths decide. It takes $O(|A|+|B|)$ time and $O(1)$
extra space beyond the supplied prefix objects, regardless of $p,q$.
Empty strings or zero exponents are handled before any modular index.
This exact version is convenient when only two powers need comparison.
\Needspace{6\baselineskip}
\lstinputlisting[style=C++,showstringspaces=false]{../strings/HashRecipes.hpp}

\subsection{Set and Multiset Fingerprints}
\textbf{Ignore order, preserve the right information.}
An integer key receives the same pair of seeded 64-bit mixer values
throughout the process. For a \emph{multiset}, add those values once
per occurrence, using intentional unsigned wraparound modulo
$2^{64}$. Removing one occurrence subtracts its pair.
Store the total count too, so different cardinalities cannot compare
equal. \texttt{MultisetHash} supports insertion, valid deletion,
multiset sum with \texttt{a+b}, and subtraction of a known
submultiset with \texttt{a-b}, all in $O(1)$.
These operations combine multisets with multiplicity, not distinct-set
union. The caller must ensure that deleted occurrences exist.

\textbf{Range composition and anagrams.}
\texttt{MultisetPrefix} builds cumulative multiset hashes in $O(n)$
time and space. \texttt{get(l,r)} subtracts the prefix before $l$,
answering range-multiset equality queries in $O(1)$.
It accepts either a \texttt{vector<long long>} or a byte string.
Signed integer keys are converted to their unique unsigned bit
patterns. In \texttt{"abccba"}, ranges $[0,3)$ and $[3,6)$ both
contain one each of \texttt{a,b,c}, although their string hashes
usually differ. For a small fixed alphabet, prefix frequency counts
are an exact alternative costing $O(\sigma)$ per query.

\textbf{A true set needs membership tracking.}
For distinct present keys, XOR their fingerprints. Insertion toggles
only when a key was absent; deletion toggles only when it was present.
\texttt{SetHash} maintains this membership in an unordered set,
so operations take expected $O(1)$ time and $O(n)$ memory overall.
The hash-table key function is randomized too.
Raw prefix XOR of an array with duplicates records \emph{odd
multiplicity}, not its set of distinct values: two equal copies cancel.
For example, \texttt{[a,a]} and \texttt{[b,b]} would both have XOR
zero. Use the sum-based prefix class for multiset/anagram queries.
Set union with overlapping keys also needs membership accounting;
XOR alone represents symmetric difference.
\Needspace{6\baselineskip}
\lstinputlisting[style=C++,showstringspaces=false]{../strings/CollectionHash.hpp}

\subsection{Zobrist Hashing for Incremental States}
Represent a state as fixed positions with values in $[0,k)$;
$-1$ means empty. A \texttt{ZobristContext(n,k)} assigns two random
words to every (position,value) pair, including emptiness.
The state hash is the XOR of all positions' selected tokens.
To replace a value, XOR out its old token and XOR in its new token.
\texttt{set(position,value)} takes $O(1)$; setting the same value
again has no effect. A move from one slot to another is two calls,
and undoing the calls restores the hash exactly.

The context costs $O(nk)$ time/space to build; a state costs $O(n)$
time/space and shares that context. Only compare states using the
\emph{same context}. Keep it alive and unchanged while any state
uses it; constructing a state directly from a temporary context is
explicitly disabled. Copies of a state retain the same context.
\texttt{key()} gives a pair of words suitable for a visited-state
set or transposition table. Include every component that affects
the actual state, such as whose turn it is, as another modeled slot.
Matching fingerprints can still collide; a full-state check is the
exact alternative for visited-state deduplication.
\Needspace{6\baselineskip}
\lstinputlisting[style=C++,showstringspaces=false]{../strings/ZobristHash.hpp}

\subsection{Dynamic String Hash: Replacements and Palindromes}
A segment-tree node stores forward and reverse \texttt{Hash} values.
If consecutive segments are $a,c$, their forward hash is $a*c$ and
their reverse hash is $\operatorname{rev}(c)*\operatorname{rev}(a)$.
An empty segment is the identity. This merge is associative but not
commutative, so range queries keep separate left and right accumulators
and merge them in their original order.

\texttt{DynamicStringHash(s)} builds in $O(n)$ time and space.
\texttt{set(i,c)} replaces one byte in $O(\log(n+1))$.
\texttt{get(l,r)} and \texttt{reversed(l,r)} return forward/reverse
hashes in $O(\log(n+1))$; \texttt{palindrome(l,r)} compares them.
An empty range is a palindrome. These hashes share the static
\texttt{PrefixHash} parameters, so static and dynamic results may be
compared directly. This class supports replacements, not insertion,
deletion or rearrangement of positions. A reverse \emph{hash} does
not mean that the tree implements a range-reversal update.
Each \texttt{Hash} stores six words and a length; the tree holds two
hashes per node, so account for its memory constant.
\Needspace{6\baselineskip}
\lstinputlisting[style=C++,showstringspaces=false]{../strings/DynamicStringHash.hpp}

\subsection{Rotation-Invariant Matrix Fingerprint}
\textbf{Purpose.} This is the lecture's specialized matrix-trace idea,
implemented as a candidate filter with an exact verification method.
A rotation changes \texttt{"abcab"} into \texttt{"cabab"}; it does
not allow arbitrary reordering of the letters.

Assign every byte an independently generated dense invertible
$3\times3$ matrix over the prime field. Generate again if inversion
fails. A string is represented by its ordered matrix product $P$.
For a split $s=uv$ with products $U,V$, its rotation $vu$ has product
$VU=U^{-1}(UV)U$. Thus rotations have the same
$\operatorname{tr}(P)$ and $\operatorname{tr}(P^2)$.
The code includes the length in its fingerprint as well.

\textbf{Range extraction.} Store ordered prefix products $A_i$ and
their inverses $I_i$. Append a byte matrix $R$ by
$A_{i+1}=A_iR$ and $I_{i+1}=R^{-1}I_i$; note the reversed order
in the inverse. The product for $[l,r)$ is $I_lA_r$.
Each query performs a constant number of $3\times3$ multiplications.
\texttt{fingerprint(l,r)} returns
$(r-l,\operatorname{tr}P,\operatorname{tr}(P^2))$ in $O(1)$ for
this fixed dimension, after $O(n)$ prefix preprocessing.
The prefix arrays alone occupy about $144(n+1)$ bytes: this is a
large constant, and two million-character objects need roughly
288 MB before their text and other allocations.

\textbf{What equality guarantees.} Unequal fingerprints rule out
rotation equivalence. Equal fingerprints are only a necessary
condition. The two traces depend on the same random matrices;
calling them two independent 61-bit hashes is unjustified.
No universal collision bound for all nonrotation word pairs is
claimed here. In particular, the nonzero-polynomial condition needed
for a Schwartz--Zippel bound has not been established for every pair.
This distinction also avoids importing the lecture's heuristic
multi-trace estimates as guarantees.

\texttt{same\_rotation(l,r,other,i,j)} first applies the fingerprint
filter, then verifies a matching candidate with KMP: search one
length-$L$ string among the $2L-1$ characters of the other's doubled
cyclic sequence. It uses $O(L)$ time and space for the verification,
and returns an \emph{exact} boolean. Matching/actually equivalent
ranges therefore take linear time; this is not an exact $O(1)$
rotation-equivalence data structure. \texttt{exact\_rotation(a,b)}
is available separately and does not require matrix preprocessing.
\Needspace{6\baselineskip}
\lstinputlisting[style=C++,showstringspaces=false]{../strings/RotationTraceHash.hpp}

\subsection{Hash Toolkit: Runnable Examples}
Build \texttt{strings/HashExamples.cpp} with C++17.
All comments give the expected output for the shown inputs.
Keep the original objects if hash matches need exact verification.
\Needspace{6\baselineskip}
\lstinputlisting[style=C++,showstringspaces=false]{../strings/HashExamples.cpp}

\begin{minipage}{\linewidth}
\textbf{Sources.} Topics inspired by
\href{https://aapc.org.ar/wp-content/uploads/2026/07/2026-07-27-lautaro-lasorsa-hashing.pdf}{Lautaro Lasorsa, Hashing (2026)}.
The implementations and examples here are independently written;
the matrix API deliberately separates its fingerprint from exact
verification. Further reading:
\href{https://cp-algorithms.com/string/string-hashing.html}{CP-Algorithms:
String Hashing}, and Fine--Wilf periodicity in
\href{https://www.isa-afp.org/browser_info/current/AFP/Combinatorics_Words/outline.pdf}{Combinatorics
on Words, Periodicity Lemma}.
\par
\end{minipage}
\par
\endgroup
